\documentclass{report}
\usepackage{amsmath,amssymb}
\setlength{\parindent}{0mm}
\setlength{\parskip}{1em}
\begin{document}
\begin{center}
\rule{6in}{1pt} \
{\large Paul Eller \\
{\bf Non-blocking conjugate gradient methods for extreme scale computing}}

201 N Goodwin Ave \\ Urbana \\ IL 61801
\\
{\tt eller3@illinois.edu}\\
William Gropp\end{center}

Many scientific and engineering applications use Krylov subspace methods
to solve large systems of linear equations. For extreme scale parallel
computing systems, the dot products in these methods (implemented using
allreduce operations in MPI) can limit performance because they are a
synchronization point or barrier. Therefore we seek to develop Krylov
subspace methods that avoid blocking allreduce operations and provide
greater parallel efficiency.

We present a rearranged preconditioned conjugate gradient method that
overlaps communication in the dot products and computation to improve
performance by hiding the cost of allreduces. These rearranged methods
optimize memory access patterns to limit the cost of using additional
vector operations that are introduced as part of the rearrangement of the
algorithm. These methods provide the ability to hide message latency by
overlapping communication with more time consuming computations so that
messages are received prior to their use. This can allow the methods to
avoid blocking due to waiting for messages and maintain better
performance despite imbalances in communication time on different nodes
caused by the network.

Performance models are developed to allow us to better understand the
best case theoretical performance of these methods with different input
parameters. We measured the performance for this method, a pipelined
conjugate gradient method, a single allreduce conjugate gradient method,
and the standard conjugate gradient method on Blue Waters in order to
improve and validate both our performance models and our implementations.


\end{document}
